\documentstyle[%


righttag%


,11pt%


,amssymb%


,russian%               remove in Eng.


,izvru%                 Set izven% in Eng.


]{amsart}





%


\newcommand{\cF}{\cal F}


\newcommand{\cT}{\cal T}


\newcommand{\End}{\operatorname{End}}


\newcommand{\tts}[1]{}





\theoremstyle{definition}


\theorembodyfont{\rm}


\newtheorem{notenonum}{\hskip\parindent Замечание}


\renewcommand{\thenotenonum}{}





%\hoffset=-15mm%                  For Eng.


\setcounter{page}{77}


\god{2001}


\nomer{3~(466)} %                        Remove in Eng.


%\nomer{Vol.\,45, No.\,3}%               Set in Eng.


%\str{\pageref{begin}--\pageref{end}}%  Set in Eng.


\udk{514.763}





\author{\it А.\,Магден, А.А.\,Салимов}


% NB:   ^^ remove in Eng.


\title{Горизонтальные лифты тензорных полей\\


на сечения тензорного расслоения}





\date{}


%\pagestyle{myheadings}%         Set in Eng.


\pagestyle{plain} %              Remove in Eng.


\begin{document}


%\thispagestyle{plain}%          Set in Eng.


\maketitle


\label{begin}





\section{Введение}





Пусть $M_n$ есть $n$-мерное дифференцируемое многообразие класса $C^\infty$,


и пусть $T^p_q(M_n)$, $p+q>0$, есть расслоение тензоров типа $(p,q)$ на $M_n$.


Цель данной статьи --- изучение горизонтальных лифтов тензорных


полей с многообразия $M_n$ на его тензорное расслоение $T^p_q (M_n)$,


рассматриваемых вдоль сечений этого расслоения. 


Применяемый нами метод с использованием оператора Вишневского 


является новым.





Введем обозначения.





1. $\pi : T^p_q(M_n) \to M_n$ ---


проекция пространства $T^p_q(M_n)$ на $M_n$.





2. Область изменения индексов $i$, $j$, $k$ --- $1, \dotsc, n$,


индексов $\ov i$, $\ov j$, $\ov k$ --- $n+1, \dotsc, n + n^{p+q}$, 


индексов $I = (i, \ov i)$, $J = (j, \ov j)$, $K = (k, \ov k)$ --- 


$1, \dotsc, n + n^{p+q}$.





3. $\cF(M_n)$ --- кольцо вещественнозначных функций класса $C^\infty$ 


на $M_n$.\; $\cT^p_q(M_n)$ --- бесконечномерное векторное пространство над 


$\Bbb R$ тензорных полей типа $(p,q)$, имеющих класс гладкости $C^\infty$.


Мы также будем рассматривать $\cT^p_q(M_n)$ как модуль над кольцом


$\cF(M_n)$. 





4. Векторные поля на $M_n$ будем обозначать через $V, W, \dotsc$ 


Тензорное поле типа $(1,1)$ будем обозначать через $\varphi$.





\section{Горизонтальные лифты векторных полей на сечения}





Обозначим через $x^j$ локальные координаты, заданные в окрестности 


$U \subset M_n$, и пусть 


$x^{\ov j} \overset{\text{def}}{=} t^{i_1 \dotsc i_p}_{j_1 \dotsc j_q}$ ---


индуцированные координаты на $\pi^{-1}(U) \subset T^p_q(M_n)$.





Тензорное поле $\alpha \in \cT^q_p(M_n)$


с помощью свертки определяет функцию на


$T^p_q(M_n)$, которую будем обозначать через $i\alpha$. 


В локальных координатах, если 


$\alpha = \alpha_{i_1 \dotsc i_p}^{j_1 \dotsc j_q} 


\partial_{j_1} \otimes \dotsb \otimes \partial_{j_q}


\otimes dx^{i_1} \otimes \dotsb\otimes dx^{i_p}$,


то 


$i\alpha (t) = \alpha_{i_1 \dotsc i_p}^{j_1 \dotsc j_q}


t^{i_1 \dotsc i_p}_{j_1 \dotsc j_q}$


для любого $t \in \cT^p_q(M_n)$.





Пусть $A \in \cT^p_q(M_n)$. Зададим вертикальный лифт 


${}^V A \in \cT^1_0(T^p_q(M_n))$ поля $A$ на $T^p_q(M_n)$ (см. [1])


условием 


${}^V A (i\alpha) = \alpha(A) \circ \pi = {}^V(\alpha(A))$, где


${}^V(\alpha(A))$ --- вертикальный лифт функции $\alpha(A) \in \cF(M_n)$. 


В карте $(x^j, x^{\ov j})$ вертикальный 


лифт ${}^V A$ поля $A$ на $T^p_q(M_n)$ имеет координаты


$$


{}^V A =


\left( \begin{array}{c} {}^V A^j \\ {}^V A^{\ov j} \end{array} \right)=


\left( \begin{array}{c} 0 \\ A^{i_1 \dotsc i_p}_{j_1 \dotsc j_q}


\end{array} \right).


$$





Пусть на $M_n$ задана аффинная связность $\nabla$ без кручения. 


Пусть $\nabla_V$ --- ковариантное дифференцирование в направлении 


$V \in \cT^1_0(M_n)$. Зададим горизонтальный лифт 


${}^H V = \ov \nabla_X$ векторного поля $V$ на 


$T^p_q(M_n)$ [1] условием


${}^H V (i \alpha) = i(\nabla_V \alpha)$,


$\alpha \in \cT^q_p(M_n)$. В карте $(x^k, x^{\ov k})$ 


компоненты ${}^H V$ имеют вид


$$


{}^H V^k = V^k, \quad 


{}^H V^{\ov k} = 


V^m\bigg(\,


\sum_{\mu=1}^q\Gamma^s_{m k_\mu}t^{l_1\dotsc l_p}_{k_1\dotsc s \dotsc k_q}-


\sum_{\lambda=1}^p \Gamma^{l_\lambda}_{m s}


t^{l_1 \dotsc s \dotsc l_p}_{k_1 \dotsc k_q}\bigg),


$$


где $\Gamma^k_{ij}$ --- компоненты связности $\nabla$ относительно локальных 


координат на $U \subset M_n$.





Пусть задано тензорное поле $\xi \in \cT^p_q(M_n)$.


Тогда соответствие $X \to \xi_X$,


где $\xi_X$ есть значение $\xi$ в $X \in M_n$, определяет отображение 


$\sigma_\xi : M_n \to T^p_q(M_n)$ такое, что


$\pi \circ \sigma_\xi = \text{id}_{M_n}$, и 


$n$-мерное подмногообразие $\sigma_\xi(M_n)$ называется сечением, заданным с 


помощью $\xi$. Если тензорное поле $\xi$ имеет локальные компоненты 


$\xi_{k_1 \dotsc k_q}^{l_1 \dotsc l_p}(x^k)$,


то сечение $\sigma_\xi(M_n)$ локально задается 


уравнениями


$x^k = x^k$, $x^{\ov k} = \xi_{k_1 \dotsc k_q}^{l_1 \dotsc l_p}(x^k)$ 


в системе координат $(x^k, x^{\ov k})$ в $T^p_q(M_n)$. Ясно, что $n$ 


векторов $B_j$, касательных к $\sigma_\xi(M_n)$, имеют координаты


$$


(B^K_j) = \left( \frac{\partial x^K}{\partial x^j} \right) = 


\left( 


\begin{array}{l} \delta^k_j \\


\partial_j \xi_{k_1 \dotsc k_q}^{l_1 \dotsc l_p}\end{array}\right) 


$$


относительно натурального репера $(\partial_k, \partial_{\ov k})$


в $T^p_q(M_n)$.





С другой стороны, слой локально задается уравнениями $x^k=\text{const}$,


$t_{k_1\dotsc k_q}^{l_1\dotsc l_p}=u_{k_1 \dotsc k_q}^{l_1 \dotsc l_p}$,


где $u_{k_1 \dotsc k_q}^{l_1 \dotsc l_p}$ --- параметры.


Поэтому, дифференцируя по 


$x^{\ov j} = t_{j_1 \dotsc j_q}^{i_1 \dotsc i_p}$,


получаем координаты


$$


(C^K_{\ov j}) = \left( \frac{\partial x^K}{\partial x^{\ov j}} \right) = 


\left( 


\begin{array}{l} 0 \\ \delta^{l_1}_{i_1} \cdots \delta^{l_p}_{i_p}


\delta^{j_1}_{k_1} \cdots \delta^{j_q}_{k_q}


\end{array}


\right) 


$$


$n^{p+q}$ векторов $C_{\ov j}$, касательных к слою,


относительно натурального репера 


$\{\partial_k, \partial_{\ov k}\}$ в $T^p_q(M_n)$.





Теперь можно взять $n+n^{p+q}$ векторных полей $B_j$, $C_{\ov j}$ 


вдоль сечения $\sigma_\xi(M_n)$. Они образуют локальное поле репера, который 


называется $(B,C)$-репером на $\sigma_\xi(M_n)$ в $\pi^{-1}(U)$.


В этом репере векторное поле ${}^HV$, ограниченное на $\sigma_\xi(M_n)$,


имеет координаты 


$$


{}^HV = 


\left( 


\begin{array}{l} {}^H{\wt V}^j \\ {}^H{\wt V}^{\ov j} 


\end{array}


\right) 


=


\left( 


\begin{array}{l} V^j \\ -\nabla_V\xi_{j_1 \dotsc j_q}^{i_1 \dotsc i_p}


\end{array}


\right).


$$





\section{Горизонтальные лифты аффинорных полей на сечения}





Для $\varphi \in \cT^1_1(M_n)$ определим тензорное поле 


${}^H\varphi \in \cT^1_1(T^p_q(M_n))$, $p+q \ge 1$,


вдоль сечения $\sigma_\xi(M_n)$ условиями


$$


\begin{cases} 


{}^H\varphi({}^HV) = {}^H(\varphi(V)) & \forall V \in \cT^1_0(M_n);\\


{}^H\varphi({}^VA) = {}^V(\varphi(A)) & \forall A \in \cT^p_q(M_n),


\end{cases}


$$


где $\varphi(A) = C(\varphi \otimes A) \in \cT^p_q(M_n)$,


и назовем ${}^H\varphi$ 


горизонтальным лифтом поля $\varphi \in \cT^1_0(M_n)$ на $T^p_q(M_n)$ 


вдоль сечения $\sigma_\xi(M_n)$.





Пусть $\xi \in \cT^p_q(M_n)$. Рассмотрим оператор Вишневского 


$$


(\Phi_\varphi \xi)_{l k_1 \dotsc k_q}^{l_1 \dotsc l_p} =


\varphi^m_l \nabla_m \xi_{k_1 \dotsc k_q}^{l_1 \dotsc l_p} -


\begin{cases}


\varphi^{l_1}_m \nabla_l \xi_{k_1 \dotsc k_q}^{m l_2 \dotsc l_p}, &


p \ge 1; \\


\varphi_{k_1}^m \nabla_l \xi_{m k_2 \dotsc k_q}^{ l_1 \dotsc l_p}, &


q \ge 1.


\end{cases}


$$





\begin{notenonum}


Пусть $\varphi$ --- интегрируемая $\varphi$-структура 


на $M_n$, и $\nabla\varphi =0$.


Если $\xi \in \cT^p_q(M_n)$ --- чистое тензорное поле 


относительно $\varphi$-структуры, то уравнение


$\Phi_\varphi \xi = 0$ задает условие 


голоморфности $\xi$ ([2], с.\,184).


\end{notenonum}





Поле ${}^H\varphi$ имеет вдоль сечения $\sigma_\xi(M_n)$ компоненты


$$


\left\{ 


\begin{array}{l} 


{}^H {\wt\varphi}^k_l = \varphi^k_l, \quad {}^H {\wt\varphi}^k_{\ov l} = 0,


\quad {}^H {\wt\varphi}^{\ov k}_l = 


- (\Phi_\varphi\xi)_{l k_1 \dotsc k_q}^{l_1 \dotsc l_p},


\\[3mm]


{}^H {\wt\varphi}^{\ov k}_{\ov l} = 


\left\{ 


\begin{array}{l} 


\varphi^{l_1}_{s_1} \delta^{l_2}_{s_2} \dotsc \delta^{l_p}_{s_p}


\delta^{r_1}_{k_1} \dotsc \delta^{r_q}_{k_q},


\quad p \ge 1;


\\[2pt]


 \delta^{l_1}_{s_1} \dotsc \delta^{l_p}_{s_p}\varphi^{r_1}_{k_1}


\delta^{r_2}_{k_2} \dotsc \delta^{r_q}_{k_q},


\quad q \ge 1,


\end{array}


\right.


\end{array}


\right.


$$


относительно $(B,C)$-репера вдоль


$\sigma_\xi(M_n)$, где $\Phi_\varphi\xi$ ---


оператор Вишневского,


$$


x^{\ov l} = t^{s_1 \dotsc s_p}_{r_1 \dotsc r_q},\quad


x^{\ov k} = t^{l_1 \dotsc l_p}_{k_1 \dotsc k_q}.


$$





В натуральном репере $\{\partial_k, \partial_{\ov k}\}$, 


заданном на $\pi^{-1}(U) \subset T^p_q(M_n)$, 


поле ${}^H\varphi$ имеет вдоль сечения $\sigma_\xi(M_n)$ компоненты


\begin{equation}


\left\{


\begin{array}{l}


{}^H {\varphi}^k_l = \varphi^k_l, \quad {}^H {\varphi}^k_{\ov l} = 0,


\quad {}^H {\varphi}^{\ov k}_{\ov l} = 


\left\{ 


\begin{array}{l} 


\varphi^{l_1}_{s_1} \delta^{l_2}_{s_2} \dotsc \delta^{l_p}_{s_p}


\delta^{r_1}_{k_1} \dotsc \delta^{r_q}_{k_q},


\quad p \ge 1; \\[2mm]


\delta^{l_1}_{s_1} \dotsc \delta^{l_p}_{s_p}\varphi^{r_1}_{k_1}


\delta^{r_2}_{k_2} \dotsc \delta^{r_q}_{k_q},


\quad q \ge 1,


\end{array}


\right.


\\[2mm]


%


{}^H\varphi^{\ov k}_l = \varphi^m_l\bigg(\displaystyle


-\sum\limits_{\lambda = 1}^p 


\Gamma^{l_\lambda}_{ms}\xi_{(k)}^{l_1 \dotsc s \dotsc l_p} +


\sum\limits_{\mu = 1}^q 


\Gamma^s_{m k_\mu}\xi^{(l)}_{k_1 \dotsc s \dotsc k_q}\bigg) +


\varphi_m^{l_1}\bigg(-\sum\limits_{\mu = 1}^q 


\Gamma^s_{lk_\mu}\xi^{m l_2 \dotsc l_p}_ {k_1 \dotsc s \dotsc k_q} +


\\[2mm]


%


\hspace{4cm}\displaystyle + \sum\limits_{\lambda = 2}^p 


\Gamma^{l_\lambda}_{ls}\xi^{m l_2 \dotsc s \dotsc l_p}_{(k)} +


\Gamma^m_{ls}\xi^{s \dotsc l_p}_{(k)}\bigg), \quad p \ge 1,


\\[2mm]


%


{}^H\varphi^{\ov k}_l = \varphi^m_l\bigg(\displaystyle


-\sum\limits_{\lambda = 1}^p 


\Gamma^{l_\lambda}_{ms}\xi_{(k)}^{l_1 \dotsc s \dotsc l_p} +


\sum\limits_{\mu = 1}^q 


\Gamma^s_{m k_\mu}\xi^{(l)}_{k_1 \dotsc s \dotsc k_q}\bigg) +


\varphi^m_{k_1}\bigg(\sum\limits_{\lambda = 1}^p 


\Gamma^{l_\lambda}_{ls}


\xi^{l_1 \dotsc s \dotsc l_p}_{m k_2 \dotsc s \dotsc k_q} -\\[2mm] 


%


\hspace{4cm}\displaystyle- \sum\limits_{\mu = 2}^q 


\Gamma^s_{l k_\mu}\xi^{(l)}_{m k_2 \dotsc s \dotsc k_q} -


\Gamma^s_{lm}\xi^{(l)}_{m \dotsc k_q}\bigg), \quad q \ge 1.


\end{array}


\right.


\end{equation}





\begin{thm}


Горизонтальный лифт ${}^H : \End M_n \to \End T^p_q(M_n)$ 


вдоль сечения $\sigma_\xi(M_n)$ есть мономорфизм.


\end{thm}





Пусть ${\frak U}_m$ --- ассоциативная 


коммутативная унитальная алгебра конечной 


размерности $m$ над $\Bbb R$. Алгебраическая $\Pi$-структура на $M_n$ --- 


это набор $\Pi = \{ \underset{\alpha}{\varphi} \}$, 


$\alpha = 1, \dotsc, m$, тензорных полей типа $(1,1)$ такой, что 


$\underset{\alpha}{\varphi} \circ \underset{\beta}{\varphi} = 


C_{\alpha\beta}^\gamma \underset{\gamma}{\varphi}$,


где $C_{\alpha\beta}^\gamma$


--- структурные константы алгебры ${\frak U}_m$.





\begin{thm}


Если $\Pi = \{ \underset{\alpha}{\varphi} \}$,


$\alpha=1,\dotsc,m$,


определяет алгебраическую 


$\Pi$-структуру на $M_n$, то набор


${}^H \Pi = \{ {}^H \underset{\alpha}{\varphi} \}$


также определяет алгебраическую


$\Pi$-структуру вдоль сечения $\sigma_\xi(M_n)$.


\end{thm}


 


\section{Новый класс квази-${\frak U}_m$-голоморфных тензорных полей}





Пусть $M_n$ и $N_m$ --- многообразия, наделенные алгебраическими структурами


$\Pi = \{ \underset{\alpha}{\varphi} \}$ и


$\wt \Pi = \{ \underset{\alpha}{\psi} \}$, $\alpha=1,\dotsc,m$,


соответственно. 


Дифференцируемое отображение $f : M_n \to N_m$ называется 


квази-${\frak U}$-голоморфным относительно $(\Pi, \wt \Pi)$ [3], 


если в каждой точке $P \in M_n$


$$


df_P \circ \underset{\alpha}{\varphi}{}_P =


\underset{\alpha}{\psi}{}_{f(P)} \circ df_P,


\quad \alpha = 1, \dotsc, m.


$$


Теперь пусть $N_k = T^p_q(M_n)$, и рассмотрим отображение 


$\sigma^\Pi_\xi : M_n \to T^p_q(M_n)$,


определяемое чистым относительно $\Pi$-структуры тензорным полем 


$\xi \in \cT^p_q(M_n)$. 


Если $\wt \Pi = \{ \underset{\alpha}{\psi} \}$ ---


алгебраическая ${}^H \Pi$-структура, 


определенная в п.\,3, то 


чистое сечение $\sigma^\Pi_\xi : M_n \to T^p_q(M_n)$ является 


квази-${\frak U}$-голоморфным относительно


$(\Pi, {}^H \Pi)$ тогда и только тогда, 


когда в локальных координатах


\begin{equation}


\underset{\alpha}{\varphi}{}^m_l\, \partial_m x^K = 


{}^H \underset{\alpha}{\varphi}{}^K_M\, \partial_l x^M,


\quad \alpha = 1, \dotsc, m,


\end{equation}


где ${}^H \underset{\alpha}{\varphi}{}^K_M$ --- компоненты поля


${}^H \underset{\alpha}{\varphi}$, заданного вдоль 


$\sigma^\Pi_\xi (M_n)$, в натуральном репере


$\{\partial_k, \partial_{\ov k} \}$ 


(см. (1)). Условие (2) приводится к виду


\begin{equation}


\underset{\alpha}{\varphi}{}^m_l


\nabla_m \xi^{l_1 \dotsc l_p}_{k_1 \dotsc k_q} -


\begin{cases}


\underset{\alpha}{\varphi}{}_m^{l_1}


\nabla_l \xi^{m l_2 \dotsc l_p}_{k_1 \dotsc k_q},& p \ge 1;


\\


\underset{\alpha}{\varphi}{}^m_{k_1}


\nabla_l \xi^{l_1 \dotsc l_p}_{m k_2 \dotsc k_q}, & q \ge 1,


\end{cases}


= 0.


\end{equation}


Таким образом, 


квази-${\frak U}$-голоморфное тензорное поле относительно


$(\Pi, {}^H \Pi)$ задается (3).





Пусть $\Pi$-структура является почти интегрируемой относительно связности


$\nabla$ без кручения, т.\,е.\


$\nabla \varphi = 0$ для любого $\varphi \in \Pi$.


Тогда получаем [4] $\Phi_\varphi \xi = {\wt\Phi}_\varphi \xi$, где


${\wt\Phi}_\varphi \xi$ --- оператор Татибаны [5]. Уравнение 


${\wt\Phi}_\varphi \xi = 0$


характеризует обычные почти голоморфные тензорные поля [5], [6]. 


Таким образом, если $\Pi$-структура почти интегрируема, то наше 


квази-${\frak U}$-голоморфное тензорное поле относительно


$(\Pi, {}^H \Pi)$ есть обычное почти голоморфное тензорное поле.


Вообще говоря, квази-${\frak U}$-голоморфное тензорное поле относительно


$(\Pi, {}^H \Pi)$, удовлетворяющее (3), не обязано удовлетворять уравнению 


${\wt\Phi}_\varphi \xi = 0$.





\begin{thebibliography}{9}





\bibitem{1}


Ledger A., Yano K. {\em Almost complex structures on tensor bundles}


// J. Diff. Geom. -- 1967. -- \No\,1. -- P.\,355--368.


\bibitem{2}


Вишневский В.В., Широков А.П., Шурыгин В.В.


{\em Пространства над алгебрами}. -- Казань:


Изд-во Казанск. ун-та, 1985. -- 262 с.


\bibitem{3}


Салимов А.А. {\em Почти $\psi$-голоморфные тензоры и их свойства} //


Докл. РАН. -- 1992. -- Т.\,324. -- \No\,3. -- С.\,533--536.


\bibitem{4}


Salimov A.A. {\em The generalized Yano--Ako operator and


complete lift of the tensor fields} //


Tensor, N.\,S. -- 1994. -- V.\,55. -- \No\,2. -- P.\,142--146.


\bibitem{5}


Tachibana S. {\em Analytic tensor and its generalization} // Tohoku Math. J.


-- 1960. -- V.\,12. -- \No\,2. -- P.\,208--221.


\bibitem{6}


Кручкович Г.И. {\em Гиперкомплексные структуры на многообразиях.} I //


Тр. семин. по векторн. и тензорн. анализу и их прилож.


-- 1972. -- Вып.\,16. -- С.\,174--201.





\end{thebibliography}





\vskip 2cm





\postup{\it Университет ``Ататюрк''}{\qquad\it Поступили}


\postup{$($\it г.\,Арзурум, Турция$)$}{\it полный текст $05.12.1997$}


\postup{}{\it краткое сообщение $23.11.1999$}


\label{end}


\end{document}


